% ============================================================
% Capítulo: Ingeniería de Requisitos (Bloque B.6bis)
% ============================================================
\chapter{Ingeniería de requisitos}
\label{cap:ingenieria-requisitos}

\section{Stakeholders}

El proyecto identifica cinco grupos de interesados reales, no
hipotéticos: los cuatro roles de usuario final del sistema --
\textbf{Lector} (estudiante o miembro de la comunidad universitaria),
\textbf{Bibliotecario} (personal de mostrador, bajo conocimiento técnico
esperado), \textbf{Gerente} (responsable de la biblioteca, además de las
funciones de Bibliotecario puede anular multas y ver reportes) y
\textbf{Administrador} (administrador técnico/institucional del
sistema) -- más el \textbf{docente-director} (Dr.\ Gleiston Cicerón
Guerrero Ulloa, Ph.D.) como evaluador académico del PFC y fuente directa
de retroalimentación formal (ver las observaciones de Entrega~1A,
Entrega~1B y Tercera Entrega en
\texttt{docs/observaciones/OBSERVACIONES.md}). El propio equipo de
desarrollo (3 integrantes) actúa simultáneamente como responsable de
elicitación, implementación y verificación -- una restricción real de
recursos declarada explícitamente en \texttt{docs/requisitos/SRS-v1.0.0.md}
(sección 2.4) que explica, entre otras cosas, por qué ciertos requisitos
quedan con evidencia parcial en vez de simularse como completos.

\section{Técnicas de elicitación}

El proyecto combina tres fuentes de requisitos, ninguna ficticia:

\begin{itemize}[leftmargin=1.2cm]
  \item \textbf{Historias de usuario y casos de uso}, en formato
    Connextra (``Como \emph{rol}, quiero \emph{acción}, para
    \emph{beneficio}'') y Cockburn respectivamente, versionados como un
    archivo por HU/CU en \texttt{docs/requisitos/historias/} y
    \texttt{docs/requisitos/casos-de-uso/} para el módulo original del
    equipo, y consolidados en dos archivos únicos
    (\texttt{historias-usuario.md}, \texttt{casos-de-uso.md}) para las 5
    HU/CU del módulo de Préstamos/Reservas/Multas -- una inconsistencia
    real de convención entre módulos, documentada explícitamente en el
    SRS en vez de normalizada en silencio.
  \item \textbf{Hallazgos de auditoría de seguridad como fuente directa
    de requisitos no funcionales}: \texttt{REQ-NF-006} (\emph{rate
    limiting} de login) y \texttt{REQ-NF-007} (auditoría de eventos de
    autenticación) no eran requisitos preexistentes -- nacieron de un
    hallazgo real de la auditoría OWASP~A07/A09 (ausencia total de
    límite de intentos y de registro de eventos), corregido dentro del
    mismo ciclo de desarrollo.
  \item \textbf{Inspección directa del código como último recurso
    declarado}, únicamente para los módulos construidos después de la
    Tercera Entrega que no llegaron a tener HU/CU dedicada -- nunca
    presentada como si una HU real existiera. De los 46 requisitos, 10
    (\texttt{REQ-F-010}, \texttt{REQ-F-020} a \texttt{REQ-F-022},
    \texttt{REQ-F-025} a \texttt{REQ-F-028}, \texttt{REQ-F-030},
    \texttt{REQ-F-031}) no citan ninguna HU/CU en la
    matriz de trazabilidad -- los últimos 2 se agregaron al cierre de
    esta entrega, verificados directamente contra \texttt{matriz.csv} --,
    y otros 5
    (\texttt{REQ-F-017}, \texttt{REQ-F-018}, \texttt{REQ-F-019}
    parcialmente, \texttt{REQ-F-023}, \texttt{REQ-F-024}) citan un
    identificador de HU/CU que \textbf{no corresponde a ningún archivo
    real} del repositorio, verificado por búsqueda exhaustiva antes de
    escribir el SRS -- en ambos casos, el rationale de esos requisitos en
    el SRS se declara explícitamente como inferido de la implementación,
    no como un hecho documentado previamente.
\end{itemize}

\section{Los 46 requisitos: categorización MoSCoW}

La fuente única de verdad es \texttt{docs/trazabilidad/matriz.csv},
validada automáticamente en cada ejecución de CI por
\texttt{scripts/validate-traceability.sh}. Este capítulo resume su
distribución agregada; el detalle de cada requisito individual
(descripción, rationale, criterio de aceptación medible, método de
verificación) vive en \texttt{docs/requisitos/SRS-v1.0.0.md} y en la
propia matriz -- no se reproducen aquí las 46 fichas completas.

\begin{center}
\begin{tabular}{|l|c|c|}
\hline
\rowcolor{verdeuteq}
\color{white}\textbf{Dimensión} & \color{white}\textbf{Categoría} & \color{white}\textbf{Cantidad (\%)} \\
\hline
\multirow{2}{*}{Tipo} & Funcional (\texttt{REQ-F}) & 31 (67{,}4\,\%) \\
 & No funcional (\texttt{REQ-NF}) & 15 (32{,}6\,\%) \\
\hline
\multirow{4}{*}{Prioridad MoSCoW} & Must & 23 (50{,}0\,\%) \\
 & Should & 23 (50{,}0\,\%) \\
 & Could & 0 (0{,}0\,\%) \\
 & Won't & 0 (0{,}0\,\%) \\
\hline
\multirow{2}{*}{Estado (matriz)} & verificado & 23 (50{,}0\,\%) \\
 & implementado & 23 (50{,}0\,\%) \\
\hline
\textbf{Total} & & \textbf{46 (100{,}0\,\%)} \\
\hline
\end{tabular}
\end{center}

Ningún requisito usa las categorías \texttt{Could} o \texttt{Won't} de
MoSCoW -- dato real, no una omisión de esta tabla: los 46 requisitos del
sistema se consideraron, al momento de incorporarse, o bien
imprescindibles (\texttt{Must}) o bien deseables dentro del alcance
actual (\texttt{Should}); ninguno se registró formalmente como
descartado o como candidato de prioridad baja.

\section{Proceso de validación: 100\,\% de los requisitos \texttt{Must} verificados}

Al cierre de este capítulo, los 23 requisitos de prioridad \texttt{Must}
tienen estado \texttt{verificado} en la matriz -- es decir, cuentan con
prueba automatizada real y/o evidencia empírica directa contra el stack
en ejecución, no solo con código que compila. Este 100\,\% es reciente:
los dos últimos requisitos \texttt{Must} en estado \texttt{implementado}
(\texttt{REQ-F-001}, registro de usuario, y \texttt{REQ-NF-013},
prevención de inyección SQL) carecían de prueba automatizada de
regresión para parte de su criterio de aceptación -- \texttt{REQ-F-001}
solo tenía cubierto el criterio de rechazo por correo duplicado, no el
flujo exitoso ni el rechazo por contraseña corta; \texttt{REQ-NF-013}
solo tenía verificación manual puntual de la auditoría original, sin
test de regresión permanente. Se cerraron agregando 2 tests nuevos
(\texttt{AuthControllerTest.registro\_passwordCorta\_devuelve400ProblemDetail}
y
\texttt{AuthServiceTest.registroConPayloadDeInyeccionSql\_} \texttt{seGuardaComoTextoLiteralSinLanzarExcepcion},
este último una regresión permanente del payload de inyección SQL
verificado manualmente en la auditoría OWASP~A03 original) en el commit
\texttt{e1f0c25}.

\section{Estabilidad del conjunto de requisitos: \texttt{CHANGELOG-REQ.md}}

\textbf{Nota de alcance (cierre de la Entrega Final).}
\texttt{SRS-v1.0.0.md} y \texttt{CHANGELOG-REQ.md} no se actualizaron
con \texttt{REQ-F-030}/\texttt{REQ-F-031} (agregados directo en
\texttt{matriz.csv} al cierre, sin tiempo de pasar por el flujo
SRS$\to$CHANGELOG-REQ); esta sección cita fielmente lo que esos dos
archivos dicen hoy (43 requisitos, no 46) en vez de inflar su conteo sin
que el archivo fuente lo respalde -- corregir ambos archivos queda fuera
de alcance de esta tarea.

\texttt{docs/requisitos/CHANGELOG-REQ.md} compara el conjunto de
requisitos entre la Tercera Entrega
(\texttt{docs/requisitos/historico/SRS-v0.9.0-rc.md}, 30 requisitos) y
esta entrega (\texttt{SRS-v1.0.0.md}, 43 requisitos), comparando el
\textbf{cuerpo completo} de cada requisito compartido, no solo su
título -- metodología que descartó explícitamente 2 falsos positivos
(\texttt{REQ-F-016}, \texttt{REQ-NF-009}) que una primera versión del
script de comparación marcó como ``modificados'' por un artefacto de
extracción de texto (un separador Markdown capturado como parte del
cuerpo), documentado en el propio \texttt{CHANGELOG-REQ.md} en vez de
contar una modificación que en realidad no ocurrió. El resultado: 13
requisitos agregados, 2 genuinamente modificados
(\texttt{REQ-NF-012}/TLS y \texttt{REQ-NF-014}/CSP, ambos reclasificados
de ``pendiente'' a ``parcialmente implementado'' tras cerrar parte de su
alcance), 0 eliminados -- \textbf{tasa de estabilidad} $=1-(2/43)=0{,}9535\approx0{,}953$,
es decir, \textbf{95{,}3\,\%}.

\textbf{Nota de honestidad (hallazgo de este capítulo, no de
\texttt{CHANGELOG-REQ.md}).} Ese mismo archivo reporta también un
``porcentaje verificado'' de $21/43$, es decir 48{,}8\,\%, calculado en
el momento en que se escribió. Esa cifra quedó desactualizada dos veces
desde entonces: primero por el cierre descrito en la sección anterior
(commit \texttt{e1f0c25}, posterior a \texttt{CHANGELOG-REQ.md}), y
después por el crecimiento de la matriz de 43 a 46 requisitos al cierre
de esta entrega (\texttt{REQ-F-030}, \texttt{REQ-F-031}, ninguno de
prioridad \texttt{Must}). El conteo real vigente de requisitos
\texttt{verificado} es $23/46$, es decir 50{,}0\,\%, no 48{,}8\,\% ni
53{,}5\,\%. Se señala
aquí en vez de propagar silenciosamente una cifra desactualizada --
\texttt{CHANGELOG-REQ.md} en sí no se corrige en esta tarea, por estar
fuera de su alcance.

\section{Checklist de calidad de requisitos (INCOSE~v4)}

\textbf{Nota de alcance (cierre de la Entrega Final).}
\texttt{docs/checklists/incose2023-req.md} tampoco se actualizó con
\texttt{REQ-F-030}/\texttt{REQ-F-031} -- evaluarlos exige releer el
texto completo de cada uno contra las 15 características C1--C15, un
trabajo analítico que no se puede improvisar de forma fiable bajo la
presión de tiempo del cierre. Esta sección, igual que la anterior, cita
fielmente los 43 requisitos que ese checklist evaluó de verdad, no 46.

\texttt{docs/checklists/incose2023-req.md} evalúa los 43 requisitos
contra las 9 características individuales (C1--C9) del marco INCOSE~v4
descrito en la \autoref{sec:mt-incose}
y las 6 características de conjunto (C10--C15), con el mismo criterio de
honestidad metodológica del resto del proyecto: no se marca cumplimiento
por defecto.

\textbf{Resultado agregado, considerando solo las 8 características de
contenido (C1--C8)}: \textbf{20 de 43 requisitos (46{,}5\,\%)} pasan las
8 sin ninguna excepción; los 23 restantes (53{,}5\,\%) tienen al menos
una excepción documentada con evidencia concreta -- típicamente C5
(\emph{Singular}: un requisito agrupa más de una acción o regla de
negocio distinguible bajo un mismo id, 12 casos) o C2
(\emph{Appropriate}: criterio de aceptación demasiado escueto frente a
requisitos hermanos del mismo módulo, 5 casos); las de C8
(\emph{Correct}) son las menos frecuentes (4 casos) pero las más
específicas, incluyendo una que este mismo checklist encontró antes de
que se cerrara: la afirmación de \texttt{REQ-F-001}/\texttt{REQ-NF-013}
de que ciertos criterios carecían de prueba automatizada, ya resuelta en
el commit \texttt{e1f0c25} citado en la sección anterior.

\textbf{Hallazgo sistémico de C9 (\emph{Conforming}), declarado sin
ambigüedad}: ninguno de los 43 requisitos está redactado como una única
cláusula ``[condición] [sujeto] \emph{shall} [acción] [objeto]
[restricción]'' del patrón formal de ISO/IEC/IEEE~29148:2018
(\autoref{sec:mt-requisitos}). El formato
real del SRS separa \emph{Descripción} (sujeto + modal + acción +
objeto, en prosa) de \emph{Criterio de aceptación medible} (las
condiciones/restricciones, en bullets aparte) -- un híbrido
Volere/Cockburn/Gherkin, no el patrón EARS de una sola oración. El
propio checklist documenta esto como una desviación
\textbf{estructural y sistemática}, no un defecto puntual de contenido
de alguno de los 43 -- razón por la cual, si se cuentan las 9
características en vez de solo las 8 de contenido, el resultado formal
es 0 de 43 requisitos sin ninguna excepción (100\,\% con al menos una),
enteramente explicado por este hallazgo de formato, no por 43 fallas de
contenido independientes. No ocultar este resultado -- aunque a primera
vista parezca alarmante -- es exactamente el criterio de honestidad
metodológica que el propio checklist declara como su objetivo.

Sobre el conjunto (C10--C15): 4 de las 6 características de conjunto se
evalúan como cumplidas con matices (\emph{Complete}, \emph{Feasible},
\emph{Comprehensible}, \emph{Able to be validated}), respaldadas por
evidencia directa (los 46 requisitos vigentes hoy en la matriz, todos
implementados y validados en CI vía
\texttt{scripts/validate-traceability.sh} -- esta es una comprobación
mecánica de CI sobre el CSV completo, no una re-evaluación INCOSE, por
eso sí se actualiza aquí a diferencia del resto de esta sección). Las
otras 2
(\emph{Consistent}, \emph{Correct}) se marcan explícitamente con
excepción: se identificaron 2 inconsistencias reales dentro del
conjunto -- la contradicción ya autodeclarada entre
\texttt{REQ-F-001}/\texttt{REQ-F-020} sobre el estado inicial de una
cuenta tras el registro, y una nueva (encontrada por el propio
checklist, no señalada antes) entre las dos cifras de latencia citadas
por \texttt{REQ-NF-003} para la misma condición de caché frío.
